\SourceLeaf{232}%
\Needspace{9\baselineskip}
\begin{center}\textsc{Bemerkungen.}\end{center}
\begingroup\small
Bereits im Juni 1831 hatte \textsc{Wolfgang Bolyai} das »Werkchen« seines Sohnes \textsc{Johann} an \textsc{Gauss} abgeschickt, es war jedoch wegen der Choleraepidemie nicht an diesen gelangt (Briefe von \textsc{Bolyai} an \textsc{Gauss} vom 20.~Juni 1831 und 16.~Januar 1832). Der Titel des Exemplares, das sich in \textsc{Gauss}' Nachlass befindet, ist von \textsc{Wolfgang Bolyai} eigenhändig geschrieben und lautet:
\begin{center}Appendix prima.\end{center}
\noindent Scientia Spatii, a veritate aut falsitate Axiomatis XI\textsuperscript{mi} Euclidei (a priori haud unquam decidenda) independens: atque ad casum falsitatis quadratura circuli geometrica.

Auctore, Auctoris filio \textsc{Johanne Bolyai}, de eadem, Geometrarum in Exercitu Caesareo Regio Austriaco Castrensium Locumtenente Primario.

De eadem bedeutet: de Bolya; Bolya war das in der Nähe von Maros Vásárhely gelegene Stammgut der Familie Bolyai (vergl. auch S.~235, Z.~18 und Fussnote).

Die nur 26 Seiten lange klassische Abhandlung \textsc{Johann Bolyais} erschien als Anhang zu dem ersten Bande des Werkes seines Vaters:

Tentamen juventutem studiosam in elementa matheseos purae, elementaris ac sublimioris, methodo intuitiva, evidentiaque huic propria, introducendi. Maros Vásárhelyini 1832.

\noindent und führt daselbst einen etwas abweichenden Titel (vergl. den Brief von \textsc{Gauss} an \textsc{Gerling} vom 4.~Februar 1844 S.~235 dieses Bandes).

Zur Erläuterung sei noch bemerkt, dass \textsc{Johann Bolyai} das Zeichen $\TripleParallel$ für \emph{parallel} (im Sinne der Nicht-Euklidischen Geometrie) anwendet und dass er mit $\Sigma$ das System der Euklidischen, mit $S$ das System seiner absoluten Geometrie bezeichnet.

Das Original des Briefes von \textsc{Gauss} vom 6.~März 1832, das \textsc{Wolfgang} seinem Sohne \textsc{Johann} geschenkt hatte, ist verloren gegangen. In \textsc{Gauss}' Nachlass befindet sich nur eine von \textsc{Johann} angefertigte Abschrift, die dessen Vater am 26.~August 1856 an \textsc{Sartorius von Waltershausen} geschickt und dieser dem \textsc{Gauss}'schen Nachlass einverleibt hat.

Die Abschrift hat S.~222 Z.~18—19 \emph{wo folglich} statt \emph{und folglich}, S.~224 Z.~18 $\Sigma$ \emph{in} $S$ statt $\Sigma$ \emph{und} $S$. Da \textsc{Gauss} häufig als Abkürzung für \emph{und} ein undeutliches $u\partial$ schreibt, so ist ein Irrthum des Abschreibers sehr wahrscheinlich. Die unterstrichenen Zeichen $\underline{f\varphi}$ und $\underline{f(180^{0}-\varphi)}$ auf S.~222, $\underline{1}$, $\underline{2}$, $\underline{3}$, $\underline{\alpha}$, $\underline{\beta}$, $\underline{\gamma}$, $\underline{A}$, $\underline{B}$, $\underline{C}$ auf S.~223 waren im Originale mit rother Tinte geschrieben und sind in der Abschrift mit Bleistift wiedergegeben.
\begin{flushright}\textsc{Stäckel.}\end{flushright}
\par\endgroup
